\documentclass[11pt,a4paper]{report}
\usepackage{polski}
\usepackage[utf8]{inputenc}
\usepackage[margin=3.5cm]{geometry}
\usepackage{graphicx}

%% ########################################################
\begin{document}
\thispagestyle{empty}
\begin{center}
\includegraphics[width=\textwidth]{logo_AGH.jpg}\\
{\bf{\sf FACULTY OF PHYSICS AND APPLIED COMPUTER SCIENCE}}\\[5mm]
%% ======================================================
{\bf{\sf{DEPARTMENT OF APPLIED COMPUTER SCIENCE AND COMPUTER PHYSICS}}}\\[14mm]
%% KATEDRA FIZYKI CIAŁA STAŁEGO
%% KATEDRA FIZYKI MATERII Skondensowanej
%% KATEDRA FIZYKI MEDYCZNEJ I BIOFIZYKI
%% KATEDRA INFORMATYKI STOSOWANEJ I FIZYKI KOMPUTEROWEJ
%% KATEDRA ODDZIAŁYWAŃ I DETEKCJI CZĄSTEK
%% KATEDRA ZASTOSOWAŃ FIZYKI JĄDROWEJ
%% ======================================================
{\sf{\huge Master Thesis}}\\[12mm] 
%% Projekt dyplomowy = inżynierska
%% Praca dyplomowa = magisterska
%% ======================================================
{\sf{\Large Augmentation of CT medical data using selected deep models\\[2mm] %% Niepotrzebne usunąć
Augmentacja danych medycznych CT przy użyciu wybranych modeli głębokich\\[2mm] %% Jeżeli praca została napisana w języku innym niż język polski
%% ======================================================
%% W przypadku pracy napisanej
%% - po polsku: dwa tytuły,
%% - po angielsku: dwa tytuły,
%% - po hiszpańsku/niemiecku/rosyjsku/etc: trzy tytuły.
%% ======================================================
}}\\[40mm]
\end{center}
{\sf
\begin{tabular}{ll}
Author: & Mateusz Górczany\\
Field of study: &	Applied Computer Science\\
%% Fizyka Techniczna

%% Informatyka Stosowana 
%% Mikro- i Nanotechnologie w biofizyce
%% Nanoinżynieria Materiałów
%% Computer Physics
Thesis supervisor: & prof. dr hab. Eng. Tomasz Szumlak\\
\end{tabular}
}\\[10mm]
\begin{center}
{\sf Kraków, 2024}
\end{center}

%% ########################################################
\newpage
%% ########################################################
\end{document}